\documentclass[10pt,a4paper]{amsart}
\usepackage[margin=19mm]{geometry}
\usepackage{amsmath,amssymb,mathtools}
\usepackage[scr=rsfs,cal=euler]{mathalfa}
\usepackage{xfrac}
\usepackage[unicode,hidelinks,bookmarks=false]{hyperref}
\newcommand{\ie}{i.e.\ }
\newcommand{\cf}{cf.\ }
\newcommand{\resp}{respectively\ }
\newcommand{\Subsection}{\subsection}
\providecommand{\nobreakdash}{\nobreak\hbox{-}}
\setlength{\emergencystretch}{2em}
\multlinegap0pt
\usepackage{amssymb}
\usepackage{natbib}
\usepackage{xcolor}
\usepackage{multirow}
\usepackage{float}
\usepackage{booktabs}
\usepackage{subfig}
\usepackage{caption}
\newtheorem{assumption}{Assumption}
\newtheorem{theorem}{Theorem}
\newtheorem{lemma}{Lemma}
\newcommand{\Leb}{\ensuremath{\mathrm{Leb}}}
\newcommand{\Levy}{L\'{e}vy }
\newcommand{\Var}{\ensuremath{\mathrm{Var}}}
\title{Nonparametric estimation of trawl processes: Theory and applications}
\author{Orimar Sauri, Almut E. D. Veraart}
\date{Source: arXiv:2209.05894v3 (CC BY 4.0)}
\begin{document}
\maketitle

\noindent\emph{Source and licence.} Nonparametric estimation of trawl processes: Theory and applications. Version arXiv:2209.05894v3 (CC BY 4.0), distributed by arXiv under the Creative Commons Attribution 4.0 International licence, http://creativecommons.org/licenses/by/4.0/, as recorded in the arXiv licence metadata for this exact version and shown on the abstract page https://arxiv.org/abs/2209.05894v3. Full licence notice: \texttt{licenses/arxiv-2209.05894-CC-BY-4.0.txt}.\par\smallskip

\noindent\textit{Editorial scope.} Counted unit: the article's lemma lemmapproxgammahat-1 (source0.tex lines 1122--1136). The article's complete original proof of it is delivered with it (source0.tex lines 1693--1981, one \texttt{proof} environment of 289 lines, which the article places away from the statement). No mathematics is written, repaired or re-derived: the statement and the proof are the article's own lines, in the article's own notation, and no retained line is substituted or re-flowed.  Retained uncounted as the source-local context the proof needs: lines 240--242; lines 244--248; lines 255--257; lines 260--268; lines 278--279; lines 1104--1119; lines 1194--1196; lines 1271--1308; lines 1434--1436; lines 1680--1688.  Nothing was omitted from any retained span, so the package carries no omission record.  No retained line contains a citation, so the wrapper carries no bibliography and no bibliography entry.  No retained line contains a figure environment, an image-inclusion command or drawing code, so no image is delivered and the package declares no asset dependency.  Editorial formatting only, in addition: an AMS \texttt{amsart} preamble that restates the article's own declarations and macros used by the retained lines, namely the environments \texttt{assumption}, \texttt{theorem}, \texttt{lemma} on separate counters, together with the standard packages those lines need.  The retained environments print in this document as Assumption 1, Theorem 1, Lemma 1 in order of appearance, whereas the article prints the same items with the numbering of its own full document.  The complete native source of the article as distributed in the arXiv e-print for this exact version is bundled unchanged as \texttt{sources/130628/source0.tex}.
\begin{assumption}\label{as:basicsamplingscheme} We assume that, as $n\uparrow\infty$,
	$\Delta_{n}\downarrow0$ and $n\Delta_{n}\rightarrow+\infty.$ 
\end{assumption}

\begin{theorem}[Consistency]\label{propconsistency}
	Let Assumption \ref{as:basicsamplingscheme} hold and suppose that $\Var(L^{\prime})=1$ and $\mathbb{E}(\mid L^{\prime}\mid^{4})<\infty$.
	Then, for all $t\geq 0$, 
	$\hat a(t) \stackrel{\mathbb{P}}{\to} a(t)$.
\end{theorem}

\begin{assumption}\label{as:samplingscheme} We assume that, as $n\uparrow\infty$,
	$\Delta_{n}\downarrow0$ and $n\Delta_{n}\rightarrow+\infty.$ Furthermore,
	$n\Delta_{n}^{3}\rightarrow\mu\in[0,\infty]$. \end{assumption}

\begin{assumption}\label{as:trawl} The trawl function admits the
	representation 
	$
	a(s)=\int_{s}^{\infty}\phi(y)dy$, for $s\geq0$, 
	where $\phi$ is a strictly positive function that is continuous on $[0,\infty)$ and such
	that $\phi(s)=\mathrm{O}(s^{-\alpha-1})$ as $s\rightarrow+\infty$
	for some $\alpha>1$.
	
\end{assumption}

\begin{assumption}\label{as:seed} $	\Var(L^{\prime})=1$
	and for some $p_{0}>4$ we have that $\mathbb{E}(\mid L^{\prime}\mid^{p_{0}})<\infty$. \end{assumption} 

\begin{lemma}\label{lemmamoments}If $L$ is a centered homogeneous
	\Levy basis, then 
	\begin{align*}
		\mathbb{E}(L(A)^{2}) & =\Leb(A)(\sigma^{2}+\int x^{2}\nu(dx)).\\
		\mathbb{E}(L(A)^{3}) & =\Leb(A)\int x^{3}\nu(dx).\\
		\mathbb{E}(L(A)^{4}) & =\Leb(A)\left\{ c_{4}(L')+3\Leb(A)(\sigma^{2}+\int x^{2}\nu(dx))^{2}\right\} ,
	\end{align*}
    where $c_{4}(L')$ is as in Theorem  \ref{propconsistency}  in the main article. Moreover, if $A\subseteq B\subseteq C$, then 
	\[
	\mathbb{E}\left[L(A)L(B)L(C)\right]=\mathbb{E}\left[L(A)^{3}\right],
	\]
	and 
	\[
	\mathbb{E}\left[L(A)^{2}L(B)L(C)\right]=\mathbb{E}\left[L(A)^{4}\right]+\mathbb{E}\left[L(A)^{2}\right]\mathbb{E}\left[L(B\backslash A)^{2}\right].
	\]
\end{lemma}

	\begin{equation}
		-\delta_{k}\tilde{X}=  \chi_{k}-\varrho_{k}; \,\,\, \tilde{X}_{(k-l_{n})\Delta_{n}}=  \beta_{k,j}^{(1)}+\beta_{k,j}^{(2)}+\sum_{m=j+1}^{n}\beta_{k,m}^{(1)}+\beta_{k}^{(3)},\label{decXbetas}
	\end{equation}

	Let us start by assuming that $t>0$ and observe that in this situation
	(recall that $L$ is assumed to be centred) 
	\[
	\mathcal{E}_{t}=-\frac{1}{\sqrt{n\Delta_{n}}}\sum_{k=l_{n}}^{n-1}\left\{ X_{(k-l_{n})\Delta_{n}}\delta_{k}X-\mathbb{E}\left(X_{(k-l_{n})\Delta_{n}}\delta_{k}X\right)\right\} +\mathrm{o}_{\mathbb{P}}(1),
	\]
	Furthermore, using (\ref{decXbetas}) and that
	\begin{equation}
		-\delta_{k}X=  \chi_{k}-\sum_{j=k+1}^{n}\alpha_{k+1,j}-\gamma_{k},\label{incremrel}
	\end{equation}
	with
	\begin{equation}
		\alpha_{k,j}=L(\mathcal{P}_{A}^{n}(k,j)),\,\,\,\gamma_{k}:=L(\cup_{j\geq n+1}\mathcal{P}_{A}^{n}(k+1,j)),\label{defalpha}
	\end{equation}
	we further deduce that 
	\begin{equation}
		\mathcal{E}_{t}=\frac{1}{\sqrt{n\Delta_{n}}}\sum_{j=l_{n}+1}^{n-1}\zeta_{j,n}+\sum_{q=1}^{3}r_{n,q}+\mathrm{o}_{\mathbb{P}}(1),\label{errorrepmtg0}
	\end{equation}
	where $\zeta_{j,n}=\sum_{p=1}^{5}\xi_{j,n}^{(p)}$ in which 
	\begin{equation}
		\begin{aligned}\xi_{j,n}^{(1)} & =\left[\chi_{j}\beta_{j,j}^{(1)}-\mathbb{E}(\chi_{j}\beta_{j,j}^{(1)})\right];\,\,\,\xi_{j,n}^{(2)}=\chi_{j}\beta_{j,j}^{(2)},\,\,\,\xi_{j,n}^{(3)}=\sum_{k=l_{n}}^{j-1}\chi_{k}\beta_{k,j}^{(1)};\\
			\xi_{j,n}^{(4)} & =-\sum_{k=l_{n}}^{j-1}\alpha_{k+1,j}\beta_{k,j}^{(2)};\,\,\,\xi_{j,n}^{(5)}=-\sum_{m=l_{n}+1}^{j}\sum_{k=l_{n}}^{m-1}\alpha_{k+1,m}\beta_{k,j}^{(1)}=-\sum_{k=l_{n}}^{j-1}\sum_{m=k+1}^{j}\alpha_{k+1,m}\beta_{k,j}^{(1)}
		\end{aligned}
		\label{errorstpos}
	\end{equation}
	Furthermore, 
\begin{equation}\label{residuals_def}
\begin{aligned}r_{n,1}= & \frac{1}{\sqrt{n\Delta_{n}}}\sum_{k=l_{n}}^{n-1}\chi_{k}\beta_{k}^{(3)},\,\,r_{n,2}=-\frac{1}{\sqrt{n\Delta_{n}}}\sum_{k=l_{n}}^{n-1}X_{(k-l_{n})\Delta_{n}}\gamma_{k}\\
r_{n,3}= & -\frac{1}{\sqrt{n\Delta_{n}}}\sum_{k=l_{n}}^{n-1}\beta_{k}^{(3)}L(\cup_{j=k+1}^{n}\mathcal{P}_{A}^{n}(k+1,j)).
\end{aligned}
\end{equation}    
	We proceed now to show that $r_{n,q}$ are asymptotically negligible
	for $q=1,2,3$. Observe first that since $(\chi_{k})_{k=l_{n},\ldots,n-1}$
	are i.i.d., centered, and independent of $(\beta_{k}^{(3)})_{k=l_{n},\ldots,n-1}$,
	then the array $(\chi_{k}\beta_{k}^{(3)})_{k=l_{n},\ldots,n-1}$ is
	a martingale difference w.r.t. the filtration 
	\[
	\mathscr{G}_{j}^{n}\lor\sigma(\{\beta_{k}^{(3)}:k=l_{n},\ldots,n-1\}),\,\,j=l_{n},\ldots,n-1.
	\]

 \begin{equation}
\mathbb{E}(\rvert L(B)\rvert^{r})\leq C(\mathrm{Leb}(B)^{r/2}\lor\mathrm{Leb}(B)),\,\,\forall\,\ensuremath{B\in\mathcal{B}(\mathbb{R}^{2}),}\label{momentineq}
\end{equation}  

\begin{lemma}\label{lemmaavar}Let Assumptions \ref{as:samplingscheme},
	\ref{as:trawl} and \ref{as:seed} in  the main article hold. 
	Then, we have that 
	\[
	\frac{1}{\Delta_{n}n}\sum_{j=l_{n}+1}^{n-1}\mathbb{E}(\zeta_{j,n}^{2}\mid\mathscr{G}_{j-1}^{n})=\frac{1}{\Delta_{n}n}\sum_{j=l_{n}+1}^{n-1}\mathbb{E}(\zeta_{j,n}^{2})+\mathrm{o}_{\mathbb{P}}(1),
	\]
	where $\zeta_{j,n}=\sum_{p=1}^{5}\xi_{j,n}^{(p)}$ and $\xi_{j,n}^{(p)}$ are defined as in (\ref{errorstpos}). 
	
\end{lemma}

\begin{lemma}\label{lemmapproxgammahat-1}Let $(\Omega,\mathscr{F},\mathbb{P})$
	be a complete probability space endowed with a sequence of discrete-time
	filtrations $(\mathscr{H}_{i}^{n})_{i\geq1}$. Suppose that $(D_{i}^{n}:i,n\geq1)$
	is a square-integrable $(\mathscr{H}_{i}^{n})_{i\geq1}$-martingale
	difference and that $N_{n}$ is an $\mathbb{N}$-valued $(\mathscr{H}_{i}^{n})_{i\geq1}$-stopping
	time increasing to $+\infty$. Then 
	\[
	\sum_{i=1}^{N_{n}}D_{i}^{n}\overset{\mathbb{P}}{\rightarrow}0,
	\]
	whenever 
	\[
	\sum_{i=1}^{N_{n}}\mathbb{E}\left[(D_{i}^{n})^{2}\mid\mathscr{H}_{i-1}^{n}\right]\overset{\mathbb{P}}{\rightarrow}0.
	\]
	
\end{lemma}

\begin{proof}[Proof of Lemma \ref{lemmaavar}] Let us start by choosing $n_{0}$ such that $l_{n}=\lfloor t/\Delta_{n}\rfloor>0$, for all $n\geq n_{0}$. As we mentioned above, most of the proof consists of writing the error terms as the sum of martingale differences and then applying Lemma \ref{lemmapproxgammahat-1}.
	Denote by $\delta_{k,j}=L(\cup_{m=k+1}^{j-1}\mathcal{P}_{A}^{n}(k+1,m))$
	and note that in this situation $
\mathbb{E}(\zeta_{j,n}^{2}\mid\mathscr{G}_{j-1}^{n})=\sum_{p,q=1}^{5}\mathbb{E}(\xi_{j,n}^{(p)}\xi_{j,n}^{(q)}\mid\mathscr{G}_{j-1}^{n})$, with a total of $15$ different summands, each of which equals (thanks
	to Lemma \ref{lemmamoments}) to 
	\begin{align}
		\begin{aligned}\mathbb{E}((\xi_{j,n}^{(1)})^{2}\mid\mathscr{G}_{j-1}^{n})= & \mathbb{E}((\xi_{j,n}^{(1)})^{2});\\
			\mathbb{E}((\xi_{j,n}^{(2)})^{2}\mid\mathscr{G}_{j-1}^{n})= & \mathbb{E}((\xi_{j,n}^{(2)})^{2})+\mathrm{O}(\Delta_{n})\left((\beta_{j,j}^{(2)})^{2}-\mathbb{E}[(\beta_{k,j}^{(2)})^{2}]\right);\\
			\mathbb{E}((\xi_{j,n}^{(3)})^{2}\mid\mathscr{G}_{j-1}^{n})= & \sum_{k,k^{\prime}=l_{n}}^{j-1}\chi_{k}\chi_{k^{\prime}}\Leb(\cup_{i=0}^{k\land k^{\prime}-l_{n}}\mathcal{P}_{A}^{n}(i,j));\\
			\mathbb{E}((\xi_{j,n}^{(4)})^{2}\mid\mathscr{G}_{j-1}^{n})= & \mathbb{E}[(\xi_{j,n}^{(4)})^{2}]+\sum_{k=l_{n}}^{j-1}\left((\beta_{k,j}^{(2)})^{2}-\mathbb{E}[(\beta_{k,j}^{(2)})^{2}]\right)\Leb(\mathcal{P}_{A}^{n}(k+1,j));\\
			\mathbb{E}((\xi_{j,n}^{(5)})^{2}\mid\mathscr{G}_{j-1}^{n})= & \sum_{k=l_{n}}^{j-1}\mathbb{E}[(\beta_{k,j}^{(1)})^{2}]\mathbb{E}[\alpha_{k+1,j}^{2}]+\sum_{k,k^{\prime}=l_{n}}^{j-2}\mathbb{E}[(\beta_{k\land k^{\prime},j}^{(1)})^{2}]\delta_{k,j}\delta_{k^{\prime},j},
		\end{aligned}
		\label{eq:condvariances}
	\end{align} 
	and 
	\begin{equation}
		\begin{aligned}\mathbb{E}(\xi_{j,n}^{(1)}\xi_{j,n}^{(2)}\mid\mathscr{G}_{j-1}^{n})=\beta_{j,j}^{(2)}\mathbb{E}[(\beta_{j,j}^{(1)})^{3}]; & \,\,\,\mathbb{E}(\xi_{j,n}^{(1)}\xi_{j,n}^{(3)}\mid\mathscr{G}_{j-1}^{n})=\sum_{k=l_{n}}^{j-1}\chi_{k}\mathbb{E}[(\beta_{k,j}^{(1)})^{3}];\\
			\mathbb{E}(\xi_{j,n}^{(1)}\xi_{j,n}^{(4)}\mid\mathscr{G}_{j-1}^{n})= & -\sum_{k=l_{n}}^{j-1}\beta_{k,j}^{(2)}\mathbb{E}[\alpha_{k+1,j}^{3}];\\
			\mathbb{E}(\xi_{j,n}^{(1)}\xi_{j,n}^{(5)}\mid\mathscr{G}_{j-1}^{n})= &-\sum_{k=l_{n}}^{j-2}\delta_{k,j}\mathbb{E}[(\beta_{k,j}^{(1)})^{3}] +\mathrm{O}(\Delta_n);\\
			\mathbb{E}(\xi_{j,n}^{(2)}\xi_{j,n}^{(3)}\mid\mathscr{G}_{j-1}^{n})= & \beta_{j,j}^{(2)}\sum_{k=l_{n}}^{j-1}\chi_{k}\mathbb{E}[(\beta_{k,j}^{(1)})^{2}];\\
			\mathbb{E}(\xi_{j,n}^{(2)}\xi_{j,n}^{(4)}\mid\mathscr{G}_{j-1}^{n})= & -\beta_{j,j}^{(2)}\sum_{k=l_{n}}^{j-1}\beta_{k,j}^{(2)}\mathbb{E}[\alpha_{k+1,j}^{2}];\\
			\mathbb{E}(\xi_{j,n}^{(2)}\xi_{j,n}^{(5)}\mid\mathscr{G}_{j-1}^{n})= & -\beta_{j,j}^{(2)}\sum_{k=l_{n}}^{j-2}\delta_{k,j}\mathbb{E}[(\beta_{k,j}^{(1)})^{2}];\\
			\mathbb{E}(\xi_{j,n}^{(3)}\xi_{j,n}^{(4)}\mid\mathscr{G}_{j-1}^{n})= & -\sum_{k,k^{\prime}=l_{n}}^{j-1}\chi_{k}\beta_{k^{\prime},j}^{(2)}\mathbb{E}[\alpha_{k^{\prime}+1,j}^{2}]\mathbf{1}_{k^{\prime}\leq k-l_{n}-1};\\
			\mathbb{E}(\xi_{j,n}^{(3)}\xi_{j,n}^{(5)}\mid\mathscr{G}_{j-1}^{n})= & -\sum_{k,k^{\prime}=l_{n}}^{j-2}\chi_{k^{\prime}}\delta_{k,j}\mathbb{E}[(\beta_{k\land k^{\prime},j}^{(1)})^{2}];\\
			\mathbb{E}(\xi_{j,n}^{(4)}\xi_{j,n}^{(5)}\mid\mathscr{G}_{j-1}^{n})= & \sum_{k,k^{\prime}=l_{n}}^{j-2}\delta_{k,j}\beta_{k^{\prime},j}^{(2)}\mathbb{E}[\alpha_{k^{\prime}+1,j}^{2}]\mathbf{1}_{k^{\prime}\leq k-l_{n}-1}.
		\end{aligned}
		\label{eq:condcovar}
	\end{equation}
	uniformly on $j$. We will verify that for all $p=1,\ldots,5$ and for all $q=p,\ldots,5$
	we have that 
	\begin{equation}
		\frac{1}{\Delta_{n}n}\sum_{j=l_{n}+1}^{n-1}\mathbb{E}(\xi_{j,n}^{(p)}\xi_{j,n}^{(q)}\mid\mathscr{G}_{j-1}^{n})=\frac{1}{\Delta_{n}n}\sum_{j=l_{n}+1}^{n-1}\mathbb{E}(\xi_{j,n}^{(p)}\xi_{j,n}^{(q)})+\mathrm{o}_{\mathbb{P}}(1).\label{eq:lemmaapproxxi}
	\end{equation}
	We start analysing the case $p=q$, for $p=2,\ldots,5$ (the case
	$p=1$ follows trivially from (\ref{eq:condvariances})).
	
	\subsubsection*{Variances}
    	\begin{description}
		\item [{$p=2$:}] In view that $\beta_{j,j}^{(2)}$ is $l_{n}$ dependent
		and $\mathbb{E}\left[\left((\beta_{j,j}^{(2)})^{2}-\mathbb{E}[(\beta_{k,j}^{(2)})^{2}]\right)^{2}\right]=\mathrm{O}(1)$
		uniformly on $j$ (due to Lemma \ref{lemmamoments}), it follows from
		the Cauchy--Schwarz inequality that 
		\begin{equation}
			\frac{1}{n^{2}}\mathrm{Var}\left[\sum_{j=l_{n}+1}^{n-1}\left((\beta_{j,j}^{(2)})^{2}-\mathbb{E}[(\beta_{k,j}^{(2)})^{2}]\right)\right]=\mathrm{O}(1/n)+\mathrm{O}(l_{n}/n)\rightarrow0,\label{variancexi2}
		\end{equation}
		which is enough.
		\item [{$p=3$:}] Plainly
		\begin{align*}
			\sum_{k,k^{\prime}=l_{n}}^{j-1}\chi_{k}\chi_{k^{\prime}}\Leb(\cup_{i=0}^{k\land k^{\prime}-l_{n}}\mathcal{P}_{A}^{n}(i,j))= & \mathbb{E}((\xi_{j,n}^{(3)})^{2})+\sum_{k=l_{n}}^{j-1}[\chi_{k}^{2}-\mathbb{E}(\chi_{k}^{2})]\Leb(\cup_{i=0}^{k-l_{n}}\mathcal{P}_{A}^{n}(i,j))\\
			& +2\sum_{k=l_{n}+1}^{j-1}\chi_{k}\sum_{k^{\prime}=l_{n}}^{k-1}\chi_{k^{\prime}}\Leb(\cup_{i=0}^{k-l_{n}}\mathcal{P}_{A}^{n}(i,j)).
		\end{align*}
		Consequently,
		\begin{align*}
			\frac{1}{n\Delta_{n}}\sum_{j=l_{n}+1}^{n-1}\mathbb{E}((\xi_{j,n}^{(3)})^{2}\mid\mathscr{G}_{j-1}^{n})= & \frac{1}{n\Delta_{n}}\sum_{j=l_{n}+1}^{n-1}\mathbb{E}((\xi_{j,n}^{(3)})^{2})\\
			& +\frac{1}{n\Delta_{n}}\sum_{k=l_{n}}^{n-2}[\chi_{k}^{2}-\mathbb{E}(\chi_{k}^{2})]\Leb(\cup_{i=0}^{k-l_{n}}\cup_{j=k+1}^{n-1}\mathcal{P}_{A}^{n}(i,j))\\
			& +2\mathrm{Var}(L^{\prime})\frac{1}{n\Delta_{n}}\sum_{k=l_{n}+1}^{n-2}\chi_{k}\nu_{k}^{(1)},
		\end{align*}
		where $\nu_{k}^{(1)}:=\sum_{k^{\prime}=l_{n}}^{k-1}\chi_{k^{\prime}}\Leb(\cup_{i=0}^{k^{\prime}-l_{n}}\cup_{j=k+1}^{n-1}\mathcal{P}_{A}^{n}(i,j))$.
		Since $(\chi_{k})_{k\geq0}$ are i.i.d., $\mathbb{E}[(\chi_{k}^{2}-\mathbb{E}(\chi_{k}^{2}))^{2}]=\mathrm{O}(\Delta_{n})$
		uniformly on $k$, and $\Leb(\cup_{i=0}^{k-l_{n}}\cup_{j=k+1}^{n+l_{n}}\mathcal{P}_{A}^{n}(i,j))\leq \Leb(A)$,
		one trivially deduce that 
		\[
		\frac{1}{n\Delta_{n}}\sum_{k=l_{n}}^{n-2}[\chi_{k}^{2}-\mathbb{E}(\chi_{k}^{2})]\Leb(\cup_{i=0}^{k-l_{n}}\cup_{j=k+1}^{n-1}\mathcal{P}_{A}^{n}(i,j))=\mathrm{o}_{\mathbb{P}}(1).
		\]
		Note now that $(\chi_{k}\nu_{k}^{(1)})_{k\geq1}$ is an $(\mathscr{F}_{k}^{n})$-martingale
		difference satisfying that 
		\[
		\mathbb{E}[(\chi_{k}\nu_{k}^{(1)})^{2}]\leq C\Delta_{n}\int_{0}^{t_{k+1}}(\int_{s}^{\infty}a(r)dr)^{2}ds.
		\]
		Therefore, 
		\begin{align*}
			\frac{1}{(n\Delta_{n})^{2}}\sum_{k=l_{n}+1}^{n-2}\mathbb{E}\left[(\chi_{k}\nu_{k}^{(1)})^{2}\right] & \leq C\frac{1}{(n\Delta_{n})^{2}}\int_{0}^{n\Delta_{n}}\int_{0}^{x}(\int_{s}^{\infty}a(r)dr)^{2}dsdx\\
			& =C\int_{0}^{1}\int_{0}^{y}(\int_{un\Delta_{n}}^{\infty}a(r)dr)^{2}dudy\rightarrow0,
		\end{align*}
		where we have done the change of variables $y=x/(n\Delta_{n})$ and
		$u=s/n\Delta_{n}$ as well as applied the Dominated Convergence Theorem.
		A simple application of Lemma \ref{lemmapproxgammahat-1} gives that
		(\ref{eq:lemmaapproxxi}) is valid in this case.
		\item [{$p=4$:}] As for $p=3,$ by decomposing 
		\begin{align*}
			&\sum_{k=l_{n}}^{j-1}\left((\beta_{k,j}^{(2)})^{2}-\mathbb{E}[(\beta_{k,j}^{(2)})^{2}]\right)\Leb(\mathcal{P}_{A}^{n}(k+1,j)) \\
			& =\sum_{i=0}^{j-l_{n}-1}\sum_{m=i}^{j-1}[\alpha_{i,m}^{2}-\mathbb{E}(\alpha_{i,m}^{2})]\sum_{k=i+l_{n}}^{(j-1)\land(m+l_{n})}\Leb(\mathcal{P}_{A}^{n}(k+1,j))\\
			& +2\sum_{i=0}^{j-l_{n}-1}\sum_{m=i+1}^{j-1}\sum_{m^{\prime}=i}^{m-1}\alpha_{i,m}\alpha_{i,m^{\prime}}\sum_{k=i+l_{n}}^{(j-1)\land(m^{\prime}+l_{n})}\Leb(\mathcal{P}_{A}^{n}(k+1,j))\\
			& +2\sum_{i=1}^{j-l_{n}-1}\sum_{i^{\prime}=0}^{i-1}\sum_{m,m^{\prime}=i}^{j-1}\alpha_{i,m}\alpha_{i^{\prime},m^{\prime}}\sum_{k=i+l_{n}}^{(j-1)\land(m^{\prime}+l_{n})\land(m+l_{n})}\Leb(\mathcal{P}_{A}^{n}(k+1,j))
		\end{align*}
		we obtain that
		\begin{equation}
			\begin{aligned}\frac{1}{n\Delta_{n}}\sum_{j=l_{n}+1}^{n-1}\mathbb{E}((\xi_{j,n}^{(4)})^{2}\mid\mathscr{G}_{j-1}^{n})= & \frac{1}{\Delta_{n}n}\sum_{j=l_{n}+1}^{n-1}\mathbb{E}((\xi_{j,n}^{(4)})^{2})\\
				& +\frac{1}{n\Delta_{n}}\sum_{i=0}^{n-2-l_{n}}\sum_{m=i}^{n-2}[\alpha_{i,m}^{2}-\mathbb{E}(\alpha_{i,m}^{2})]\theta_{i,m,m^{\prime}}^{(1)}\\
				& +\frac{2}{n\Delta_{n}}\sum_{i=0}^{n-2-l_{n}}\sum_{m=i+1}^{n-2}\sum_{m^{\prime}=i}^{m-1}\alpha_{i,m}\alpha_{i,m^{\prime}}\theta_{i,m,m^{\prime}}^{(2)}\\
				& +\frac{2}{n\Delta_{n}}\sum_{i=0}^{n-2-l_{n}}\sum_{i^{\prime}=0}^{i-1}\sum_{m,m^{\prime}=i}^{n-2}\alpha_{i,m}\alpha_{i^{\prime},m^{\prime}}\theta_{i,m,m^{\prime}}^{(3)},
			\end{aligned}
			\label{eq:decompq4}
		\end{equation}
		where 
		\begin{equation}
			\begin{aligned}\theta_{i,m,m^{\prime}}^{(1)} & =\sum_{j=(i+l_{n}+1)\lor(m+1)}^{n-1}\sum_{k=i+l_{n}+1}^{j\land(m+l_{n}+1)}\Leb(\mathcal{P}_{A}^{n}(k,j));\\
				\theta_{i,m,m^{\prime}}^{(2)} & =\sum_{j=(i+l_{n}+1)\lor(m+1)}^{n-1}\sum_{k=i+l_{n}+1}^{j\land(m^{\prime}+l_{n}+1)}\Leb(\mathcal{P}_{A}^{n}(k,j));\\
				\theta_{i,m,m^{\prime}}^{(3)} & =\sum_{j=(i+l_{n}+1)\lor(m+1)\lor(m^{\prime}+1)}^{n-1}\sum_{k=i+l_{n}+1}^{j\land(m^{\prime}+l_{n}+1)\land(m+l_{n}+1)}\Leb(\mathcal{P}_{A}^{n}(k,j)).
			\end{aligned}
			\label{eq:that123def}
		\end{equation}
		Observe that all the error terms in (\ref{eq:decompq4}) are sums
		of $(\mathscr{F}_{i}^{n})$-martingale differences and for $\ell=1,2,3$
		\begin{align*}
			\left|\theta_{i,m,m^{\prime}}^{(\ell)}\right| & \leq\sum_{j=(i+l_{n}+1)\lor(m+1)}^{m+l_{n}}\sum_{k=i+l_{n}+1}^{j}\Leb(\mathcal{P}_{A}^{n}(k,j))+\sum_{j=m+l_{n}+1}^{n}\sum_{k=i+l_{n}+1}^{m+l_{n}+1}\Leb(\mathcal{P}_{A}^{n}(k,j))\\
			& \leq C\Delta_{n}l_{n}+\Leb(A)=\mathrm{O}(1).
		\end{align*}
		An application of the previous bound, Lemma \ref{lemmamoments}, 
		and the independent scattered property of $L$ gives us that 
		\begin{align*}
			\sum_{i=0}^{n-2-l_{n}}\mathbb{E}\left[\left(\sum_{m=i}^{n-2}[\alpha_{i,m}^{2}-\mathbb{E}(\alpha_{i,m}^{2})]\theta_{i,m,m^{\prime}}^{(1)}\right)^{2}\right] & \leq C\sum_{i=0}^{n}\sum_{m=i}^{n}\Leb(\mathcal{P}_{A}^{n}(i,m))\leq C\Delta_{n}n,\\
			\sum_{i=0}^{n-2-l_{n}}\mathbb{E}\left[\left(\sum_{m=i+1}^{n-2}\alpha_{i,m}\sum_{m^{\prime}=i}^{m-1}\alpha_{i,m^{\prime}}\theta_{i,m,m^{\prime}}^{(2)}\right)^{2}\right] & \leq C\sum_{i=0}^{n}\left(\sum_{m=i}^{n}\Leb(\mathcal{P}_{A}^{n}(i,m))\right)^{2}\leq C\Delta_{n}n,
		\end{align*}
		and that 
		\begin{equation}
			\begin{aligned}\sum_{i=0}^{n-2-l_{n}}\mathbb{E}\left(\sum_{i^{\prime}=0}^{i-1}\sum_{m,m^{\prime}=i}^{n-2}\alpha_{i,m}\alpha_{i^{\prime},m^{\prime}}\right)^{2}= & \sum_{i=0}^{n-2-l_{n}}\sum_{i^{\prime}=0}^{i-1}\sum_{m,m^{\prime}=i}^{n-2}\mathbb{E}[\alpha_{i,m}^{2}]\mathbb{E}[\alpha_{i^{\prime},m^{\prime}}](\theta_{i,m,m^{\prime}}^{(3)})^{2}\\
				& \leq C\sum_{i=1}^{n}\sum_{m=i}^{n}\sum_{i^{\prime}=0}^{i}\sum_{m^{\prime}=i}^{n}\Leb(\mathcal{P}_{A}^{n}(i,m))\\
				&\times\Leb(\mathcal{P}_{A}^{n}(i^{\prime},m^{\prime}))\\
				& \leq C\sum_{i=0}^{n}\sum_{m=i}^{n}\Leb(\mathcal{P}_{A}^{n}(i,m))\\
				& \leq C\Delta_{n}n.
			\end{aligned}
			\label{eq:bound2varp4}
		\end{equation}
		(\ref{eq:lemmaapproxxi}) is obtained as a simple application of the
		previous previous estimates along with Lemma  \ref{lemmamoments}.
		\item [{$p=5$:}] Since in this situation
		\begin{align*}
			\frac{1}{n\Delta_{n}}\sum_{j=l_{n}+1}^{n-1}\mathbb{E}((\xi_{j,n}^{(5)})^{2}\mid\mathscr{G}_{j-1}^{n})= & \frac{1}{n\Delta_{n}}\sum_{j=l_{n}+1}^{n-1}\mathbb{E}((\xi_{j,n}^{(5)})^{2})\\
			& +\frac{1}{n\Delta_{n}}\sum_{k=l_{n}+1}^{n-2}\sum_{m=k}^{n-2}\left(\alpha_{k,m}^{2}-\mathbb{E}[\alpha_{k,m}^{2}]\right)\theta_{k,m}^{(4)}\\
			& +\frac{2}{n\Delta_{n}}\sum_{k=l_{n}+1}^{n-3}\sum_{m=k+1}^{n-2}\sum_{m^{\prime}=k}^{m-1}\alpha_{k,m}\alpha_{k,m^{\prime}}\theta_{k,m}^{(4)}\\
			& +\frac{2}{n\Delta_{n}}\sum_{k=l_{n}+2}^{n-2}\sum_{m=k}^{n-2}\sum_{k^{\prime}=l_{n}+1}^{k-1}\sum_{m^{\prime}=k}^{n-2}\alpha_{k,m}\alpha_{k^{\prime},m^{\prime}}\theta_{k,m,m^{\prime}}^{(5)},
		\end{align*}
		with the error term being once again sums of $(\mathscr{F}_{k}^{n})$-martingale
		differences, and where 
		\begin{equation}
			\theta_{k,m}^{(4)}=\Leb(\cup_{i=0}^{k-1-l_{n}}\cup_{j=m+1}^{n-1}\mathcal{P}_{A}^{n}(i,j));\,\,\theta_{k,m,m^{\prime}}^{(5)}=\Leb(\cup_{i=0}^{k-1-l_{n}}\cup_{j=m\lor m^{\prime}+1}^{n-1}\mathcal{P}_{A}^{n}(i,j)),\label{eq:theta45def}
		\end{equation}
		are uniformly bounded by $\Leb(A)$, we can argue as in the case $p=4$ to obtain
		the desired result.
	\end{description}
	
	\subsubsection*{Covariances}
		In this part we verify that (\ref{eq:lemmaapproxxi}) holds when $p\neq q$,
	for $p=1,\ldots,4$. Below we will use the notation $v_{3}=\int x^{3}\nu(dx)$.
	\begin{description}
		\item [{$p=1$:}] The case $q=2$ can be analysed similarly to the case
		$p=q=2$. Furthermore, since
		\begin{align*}
			\frac{1}{\Delta_{n}n}\sum_{j=l_{n}+1}^{n-1}\mathbb{E}(\xi_{j,n}^{(1)}\xi_{j,n}^{(3)}\mid\mathscr{G}_{j-1}^{n})= & v_{3}\frac{1}{\Delta_{n}n}\sum_{k=l_{n}}^{n-2}\chi_{k}\theta_{k-1,k}^{(4)},\\
			\frac{1}{\Delta_{n}n}\sum_{j=l_{n}+1}^{n-1}\mathbb{E}(\xi_{j,n}^{(1)}\xi_{j,n}^{(4)}\mid\mathscr{G}_{j-1}^{n})= & -\frac{v_{3}}{n\Delta_{n}}\sum_{i=0}^{n-1-l_{n}}\sum_{m=i}^{n-2}\alpha_{i,m}\theta_{i,m,m^{\prime}}^{(1)},\\
			\frac{1}{\Delta_{n}n}\sum_{j=l_{n}+1}^{n-1}\mathbb{E}(\xi_{j,n}^{(1)}\xi_{j,n}^{(5)}\mid\mathscr{G}_{j-1}^{n})= & -\frac{v_{3}}{n\Delta_{n}}\sum_{k=l_{n}+1}^{n-2}\sum_{m=k}^{n-2}\alpha_{k,m}\theta_{k,m}^{(4)}+\mathrm{O}(\Delta_n),
		\end{align*}
		in which $\theta_{i,m,m^{\prime}}^{(1)}$ and $\theta_{k,m}^{(4)}$
		are is as in  (\ref{eq:that123def}) and (\ref{eq:theta45def}), respectively,
		we can reproduce the argument used above for $p=q$, $p=3,4,5$, to
		conclude that (\ref{eq:lemmaapproxxi}) is also satisfied in this
		situation. 
		\item [{$p=2$:}] By the Cauchy-Schwartz inequality and Assumption 
		\eqref{as:trawl}
		in the main article, 
		for all $l_{n}+1\leq j<2l_{n}$
		\[
		\mathbb{E}\left[\left|\mathbb{E}(\xi_{j,n}^{(2)}\xi_{j,n}^{(3)}\mid\mathscr{G}_{j-1}^{n})\right|\right]\leq C\mathbb{E}\left[\left(\sum_{k=l_{n}}^{j-1}\chi_{k}\Leb(\cup_{i=0}^{k-l_{n}}\mathcal{P}_{A}^{n}(i,j))\right)^{2}\right]^{1/2}\leq C\Delta_{n},
		\]
		which easily implies that
		\[
		\frac{1}{\Delta_{n}n}\sum_{j=l_{n}+1}^{2l_{n}}\mathbb{E}(\xi_{j,n}^{(2)}\xi_{j,n}^{(3)}\mid\mathscr{G}_{j-1}^{n})\overset{\mathbb{P}}{\rightarrow}0.
		\]
		For $j\geq2l_{n}+1$ we have that 
		\begin{align*}
			\mathbb{E}(\xi_{j,n}^{(2)}\xi_{j,n}^{(3)}\mid\mathscr{G}_{j-1}^{n})= & \beta_{j,j}^{(2)}\vartheta_{j}^{(1)}+\beta_{j,j}^{(2)}\vartheta_{j}^{(2)},
		\end{align*}
		where 
		\[
		\vartheta_{j}^{(1)}:=\sum_{k=l_{n}}^{j-l_{n}-1}\chi_{k}\mathbb{E}[(\beta_{k,j}^{(1)})^{2}];\,\,\,\vartheta_{j}^{(2)}:=\sum_{k=j-l_{n}}^{j-1}\chi_{k}\mathbb{E}[(\beta_{k,j}^{(1)})^{2}]
		\]
		Note that $\beta_{j,j}^{(2)}\vartheta_{j}^{(1)}$ and $\beta_{j,j}^{(2)}\vartheta_{j}^{(2)}-\mathbb{E}(\beta_{j,j}^{(2)}\vartheta_{j}^{(2)})$
		are $l_{n}$-uncorrelated. Furthermore, by Rosenthal's inequality
		and relation \eqref{momentineq}, we get that 
		\[
		\mathbb{E}[\mid\vartheta_{j}^{(v)}\mid{}^{4}]\leq C\Delta_{n}^{4},\,\,\,v=1,2,
		\]
		which in turn implies that 
		\begin{align*}
			\mathbb{E}[(\beta_{j,j}^{(2)}\vartheta_{j}^{(v)})^{2}] & \leq C\Delta_{n}^{2},\,\,\,v=1,2.
		\end{align*}
		Using this bound, we deduce as in (\ref{variancexi2}) that 
		\[
		\frac{1}{n\Delta_{n}}\sum_{j=2l_{n}+1}^{n-1}\beta_{j,j}^{(2)}\vartheta_{j}^{(1)}+\frac{1}{n\Delta_{n}}\sum_{j=2l_{n}+1}^{n-1}\left[\beta_{j,j}^{(2)}\vartheta_{j}^{(2)}-\mathbb{E}(\beta_{j,j}^{(2)}\vartheta_{j}^{(2)})\right]=\mathrm{o}_{\mathbb{P}}(1),
		\]
		which is clearly enough for (\ref{eq:lemmaapproxxi}) for the case
		$q=3$. Furthermore, since for $j\geq2l_{n}+1$
		\begin{align*}
			\mathbb{E}(\xi_{j,n}^{(2)}\xi_{j,n}^{(4)}\mid\mathscr{G}_{j-1}^{n})= & -\beta_{j,j}^{(2)}(\vartheta_{j}^{(3)}+\vartheta_{j}^{(4)}),\\
			\mathbb{E}(\xi_{j,n}^{(2)}\xi_{j,n}^{(5)}\mid\mathscr{G}_{j-1}^{n})= & -\beta_{j,j}^{(2)}(\vartheta_{j}^{(5)}+\vartheta_{j}^{(6)}+\vartheta_{j}^{(7)});
		\end{align*}
		where 
		\begin{align*}
			\vartheta_{j}^{(3)}:=\sum_{i=0}^{j-l_{n}-1}\sum_{m=i}^{j-l_{n}-1}\alpha_{i,m}\theta_{i,m,j}^{(6)}; & \,\,\vartheta_{j}^{(4)}:=\sum_{i=0}^{j-l_{n}-1}\sum_{m=j-l_{n}}^{j-1}\alpha_{i,m}\theta_{i,m,j}^{(6)};\\
			\vartheta_{j}^{(5)}:=\sum_{k=l_{n}+1}^{j-l_{n}}\sum_{m=k}^{j-l_{n}-1}\alpha_{k,m}\theta_{k,j}^{(7)}; & \,\,\vartheta_{j}^{(6)}:=\sum_{k=l_{n}+1}^{j-l_{n}}\sum_{m=j-l_{n}}^{j-1}\alpha_{k,m}\theta_{k,j}^{(7)},
		\end{align*}
		and $\vartheta_{j}^{(7)}=\sum_{k=j-l_{n}+1}^{j-1}\sum_{m=k}^{j-1}\alpha_{k,m}\theta_{k,j}^{(7)}$,
		in which
		\begin{align*}
			\theta_{i,m,j}^{(6)}\lor\theta_{i,m,j}^{(7)} & \leq C\Delta_{n}a(t_{j}-t_{m}-(l_{n}+1)\Delta_{n}),
		\end{align*}
		and $\beta_{j,j}^{(2)}\vartheta_{j}^{(l)}-\mathbb{E}(\beta_{j,j}^{(2)}\vartheta_{j}^{(v)})$
		are $l_{n}$-uncorrelated for $v=3,4,5,6$, we can repeat the preceding
		arguments to conclude that (\ref{eq:lemmaapproxxi}) is also valid
		for these cases.
		\item [{$p=3$:}] By decomposing 
		\begin{equation}
			\beta_{k^{\prime},j}^{(2)}=\beta_{k^{\prime},k}^{(1)}+\beta_{k^{\prime},k}^{(2)}+\sum_{m=k+1}^{j-1}\beta_{k^{\prime},m}^{(1)},\label{eq:decompp34}
		\end{equation}
		for $k^{\prime}\leq k\leq j-2$, we obtain, due to (\ref{eq:condcovar}),
		that 
		\begin{equation}
			\begin{aligned}\sum_{j=l_{n}+1}^{n-1}\mathbb{E}(\xi_{j,n}^{(3)}\xi_{j,n}^{(4)}\mid\mathscr{G}_{j-1}^{n})= & \sum_{j=l_{n}+1}^{n-1}\mathbb{E}(\xi_{j,n}^{(3)}\xi_{j,n}^{(4)})\\
				& +\sum_{m=2l_{n}+2}^{n-1}\sum_{k=2l_{n}+1}^{m-1}\sum_{k^{\prime}=l_{n}}^{k-l_{n}-1}\chi_{k}\beta_{k^{\prime},m}^{(1)}\theta_{m,k^{\prime}}^{(8)}\\
				& +\sum_{k=2l_{n}+1}^{n-2}\sum_{k^{\prime}=l_{n}}^{k-l_{n}-1}\chi_{k}\beta_{k^{\prime},k}^{(2)}\theta_{k,k^{\prime}}^{(8)}\\
				& +\sum_{k=2l_{n}+1}^{n-2}\sum_{k^{\prime}=l_{n}}^{k-l_{n}-1}[\chi_{k}\beta_{k^{\prime},k}^{(1)}-\mathbb{E}(\chi_{k}\beta_{k^{\prime},k}^{(1)})]\theta_{k,k^{\prime}}^{(8)},
			\end{aligned}
			\label{eq:decompp3q4}
		\end{equation}
		where $\theta_{k,k^{\prime}}^{(8)}:=\Leb(\cup_{j=k+1}^{n}\mathcal{P}_{A}^{n}(k^{\prime}+1,j)).$
		Note that once again all the error terms are sums of $(\mathscr{G}_{k})$-increment
		martingales. Due to Lemma \ref{lemmapproxgammahat-1}, we only need to verify
		that the second moment of each of these summands is of order $\mathrm{O}(\Delta_{n})$.
		For the last two sums, this follows easily by using that $\mathbb{E}[(\beta_{k^{\prime},k}^{(2)})^{2}]\leq C,$
		and
		\[
		\mathbb{E}(\chi_{k}^{2}\beta_{k^{\prime},k}^{(1)}\beta_{m,k}^{(1)})\leq C\Delta_{n}a(t_{k}-t_{k^{\prime}}\land t_{m}),\,\,\,\left|\theta_{k,k^{\prime}}^{(8)}\right|\leq C\Delta_{n}a(t_{k}-t_{k^{\prime}}).
		\]
		The required bound for the second sum in (\ref{eq:decompp3q4}) is
		a bit more involved. Let us start by defining  $\vartheta_{m}^{(8)}=\sum_{k=2l_{n}+1}^{m-1}\sum_{k^{\prime}=l_{n}}^{k-l_{n}-1}\chi_{k}\beta_{k^{\prime},m}^{(1)}\theta_{m,k^{\prime}}^{(8)}$.
		In view that $\mathbb{E}[(\beta_{q,m}^{(1)})^{2}]\leq C\Delta_{n}a(t_{m}-t_{q})$,
		the sought-after bound is obtained by the following estimates
		\begin{align*}
			\mathbb{E}[(\vartheta_{m}^{(8)})^{2}]\leq & C\Delta_{n}\sum_{k=2l_{n}+1}^{m-1}\sum_{k^{\prime}=l_{n}}^{k-l_{n}-1}\sum_{q=l_{n}}^{k^{\prime}}\mathbb{E}[(\beta_{q,m}^{(1)})^{2}]\theta_{m,k^{\prime}}^{(8)}\theta_{m,q}^{(8)}\\
			\leq & C\Delta_{n}^{2}\sum_{k=2l_{n}+1}^{m-1}\int_{l_{n}\Delta_{n}}^{t_{k+1}-l_{n}\Delta_{n}}a(t_{m}-r)\left(\int_{l_{n}\Delta_{n}}^{r}a(t_{m}-s)^{2}ds\right)dr\\
			= & C\Delta_{n}^{2}\sum_{k=2l_{n}+1}^{m-1}\int_{t_{m}-t_{k+1}+l_{n}\Delta_{n}}^{t_{m}-l_{n}\Delta_{n}}a(x)\left(\int_{x}^{t_{m}-l_{n}\Delta_{n}}a(y)^{2}dy\right)dx\\
			\leq & C\Delta_{n}^{2}\sum_{k=2l_{n}+1}^{m-1}\int_{t_{m}-t_{k+1}+l_{n}\Delta_{n}}^{t_{m}-l_{n}\Delta_{n}}a(y)^{2}dy\\
			\leq & C\Delta_{n}\int_{(2l_{n}+1)\Delta_{n}}^{t_{m}}\int_{t_{m}-u+(l_{n}-1)\Delta_{n}}^{t_{m}-l_{n}\Delta_{n}}a(y)^{2}dydu\\
			\leq & C\Delta_{n}\int_{(2l_{n}+1)\Delta_{n}}^{t_{m}-l_{n}\Delta_{n}}\int_{w}^{t_{m}-l_{n}\Delta_{n}}a(y)^{2}dydw\\
			\leq & C\Delta_{n}\int_{0}^{t_{m}}a(y)^{2}ydy\leq C\Delta_{n},
		\end{align*}
		where in the last step we used the fact that $a(s)=\mathrm{O}(s^{-\alpha})$,
		as $s\rightarrow+\infty$ for some $\alpha>1$. For $q=5$, we use
		the decomposition 
		\[
		\sum_{q=k+1}^{j-2}\chi_{q}=L(\cup_{i=0}^{k}\cup_{q=k+1}^{j-2}\mathcal{P}_{A}^{n}(i,q))+\sum_{i=k}^{j-3}\delta_{i,j-1},
		\]
		to deduce that 
		\begin{align*}
			\sum_{j=l_{n}+1}^{n-1}\mathbb{E}(\xi_{j,n}^{(3)}\xi_{j,n}^{(5)}\mid\mathscr{G}_{j-1}^{n})= & \sum_{j=l_{n}+1}^{n-1}\mathbb{E}(\xi_{j,n}^{(3)}\xi_{j,n}^{(5)})+\sum_{j=l_{n}+1}^{n-1}\vartheta_{j}^{(9)}\\
			& +\sum_{k=l_{n}}^{n-3}\sum_{m=k+1}^{n-2}\alpha_{k+1,m}\vartheta_{k,m}^{(10)},
		\end{align*}
		where $\vartheta_{k,m}^{(10)}  =\sum_{q=l_{n}}^{k}\chi_{q}\theta_{q+1,m,m}^{(5)}$, and
		\begin{align*}
			\vartheta_{j}^{(9)}= & \sum_{k=l_{n}}^{j-3}\sum_{i=k}^{j-3}\mathbb{E}[(\beta_{k,j}^{(1)})^{2}][\delta_{k,j}\delta_{k,j-1}-\mathbb{E}(\delta_{k,j}\delta_{k,j-1})]\\
			& +\sum_{k=l_{n}}^{j-3}\delta_{k,j}L(\cup_{i=0}^{k}\cup_{q=k+1}^{j-2}\mathcal{P}_{A}^{n}(i,q)).	
		\end{align*}
		Note that the sum associated with $\vartheta_{j}^{(9)}$ can be analysed
		in the same way as in the case $p=q=5$. Moreover, the last error
		term are sums of $(\mathscr{F}_{k})$-increment martingale. After
		some simple changes of variables and order of integration we further
		have that
		\begin{equation*}
			\sum_{k=l_{n}}^{n-3}\sum_{m=k+1}^{n-2}\mathbb{E}[(\alpha_{k+1,m}\vartheta_{k,m}^{(10)})^{2}]  \leq C\int_{0}^{n\Delta_{n}}\int_{0}^{r}\varGamma(s)^{2}dsdr=\mathrm{o}[(n\Delta_{n})^{2}],
		\end{equation*}
		which according to Lemma \ref{lemmaavar}, let us conclude that (\ref{eq:lemmaapproxxi})
		is also valid in this case.
		\item [{$p=4$}] Lastly, from (\ref{eq:decompp34}), we obtain easily that
		\[
		\sum_{j=l_{n}+1}^{n-1}\mathbb{E}(\xi_{j,n}^{(4)}\xi_{j,n}^{(5)}\mid\mathscr{G}_{j-1}^{n})=\sum_{k=2l_{n}+1}^{n-3}\sum_{m=k+1}^{n-2}\alpha_{k+1,m}\vartheta_{k,m}^{(11)},
		\]
		with 
		\[
		\vartheta_{k,m}^{(11)}=\sum_{q=l_{n}}^{k-l_{n}-1}\sum_{p=q-l_{n}}^{n-2}\beta_{q,p}^{(1)}\Leb(\cup_{j=m\lor p+1}^{n-2}\mathcal{P}_{A}^{n}(q+1,j)).
		\]
		Since $\alpha_{k+1,m}\vartheta_{k,m}^{(11)}$ is an $(\mathscr{F}_{k})$-increment
		martingale and for $m\geq k+1$, $\mathbb{E}[(\vartheta_{k,m}^{(11)})^{2}]$ is bounded up to a constant by
		\begin{align*} \sum_{q=l_{n}}^{k-l_{n}-1}\sum_{q^{\prime}=l_{n}}^{q}&\Leb(\cup_{i=0}^{q^{\prime}-l_{n}}\cup_{p=q-l_{n}}^{n-2}\mathcal{P}_{A}^{n}(q+1,p))\Leb(\cup_{j=m+1}^{n-2}\mathcal{P}_{A}^{n}(q+1,j))\\
			&\times\Leb(\cup_{j=m+1}^{n-2}\mathcal{P}_{A}^{n}(q^{\prime}+1,j))	\end{align*}
		which is in turn uniformly bounded, we deduce as before that (\ref{eq:lemmaapproxxi}) holds in this situation.
	\end{description}
	
\end{proof}

\end{document}
